\documentclass[11pt]{article}
\usepackage[margin=1in]{geometry}
\usepackage{natbib}
\usepackage{booktabs}
\usepackage{amsmath}
\usepackage{graphicx}
\usepackage{hyperref}
\usepackage{authblk}

\title{Thiolutin Has Complex Effects \textit{In Vivo} but Is a Direct Inhibitor of RNA Polymerase II \textit{In Vitro}}

\author[1]{Authors}
\affil[1]{Department of Biochemistry, University of Pittsburgh, Pittsburgh, PA, USA}

\date{}

\begin{document}
\maketitle

\begin{abstract}
Thiolutin is a natural product transcription inhibitor with an unresolved mode of action. Thiolutin and the related dithiolopyrrolone holomycin chelate Zn$^{2+}$, and previous studies have concluded that RNA Polymerase II (Pol~II) inhibition \textit{in vivo} is indirect. Here, we present chemicogenetic and biochemical approaches to investigate thiolutin's mode of action in \textit{Saccharomyces cerevisiae}. We identify mutants that alter sensitivity to thiolutin. We provide genetic evidence that thiolutin causes oxidation of thioredoxins \textit{in vivo} and that thiolutin both induces oxidative stress and interacts functionally with multiple metals including Mn$^{2+}$ and Cu$^{2+}$, and not just Zn$^{2+}$. Finally, we show direct inhibition of RNA polymerase~II (Pol~II) transcription initiation by thiolutin \textit{in vitro} in support of classical studies that thiolutin can directly inhibit transcription \textit{in vitro}. Inhibition requires both Mn$^{2+}$ and appropriate reduction of thiolutin as excess DTT abrogates its effects. Pause-prone, defective elongation can be observed \textit{in vitro} if inhibition is bypassed. Thiolutin effects on Pol~II occupancy \textit{in vivo} are widespread but major effects are consistent with prior observations for Tor pathway inhibition and stress induction, suggesting that thiolutin use \textit{in vivo} should be restricted to studies on its modes of action and not as an experimental tool.
\end{abstract}

\section{Introduction}

Thiolutin is a historically used transcription inhibitor with potentially multiple modes of action \citep{jimenez1973,tipper1982}. It is one of the microbially produced dithiolopyrrolone compounds that feature an intramolecular and redox-sensitive disulfide bond \citep{li2014dithiolopyrrolones,li2011holomycin}. Holomycin is another notable compound in this class. These compounds are thought to be pro-drugs whose antimicrobial action requires reduction of the intramolecular disulfide \citep{li2011holomycin}. Thiolutin has been shown to inhibit bacterial and eukaryotic transcription \textit{in vivo} and has been used to study mRNA stability in multiple species \citep{grigull2004,millan2006,parker2003,morey2013}.

However, several lines of evidence have complicated the interpretation of thiolutin's mechanism of action. The transcriptional inhibitor thiolutin blocks mRNA degradation in yeast \citep{pelechano2008}, and alterations of signaling pathways in response to chemical perturbations used to measure mRNA decay rates have been documented \citep{guan2020}. Recent studies have identified thiolutin as a zinc chelator that inhibits the Rpn11 and other JAMM metalloproteases \citep{robertson2018}, and dithiolopyrrolones have been shown to disrupt bacterial metal homeostasis \citep{chan2017}. Most recently, dithiolopyrrolones have been characterized as prochelators activated by glutathione \citep{lauinger2023}.

Despite decades of study, whether thiolutin directly inhibits eukaryotic RNA polymerase II remains contested, with recent publications concluding that \textit{in vivo} Pol~II inhibition is indirect. Here, we combine chemicogenetic screens in \textit{S. cerevisiae} with biochemical reconstitution to revisit this question.

\section{Results}

\subsection{Thiolutin or reduced thiolutin alone fail to inhibit purified Pol~II \textit{in vitro}}

Given that thiolutin and holomycin activities could be sensitive to reduction of the disulfide bond by DTT \citep{li2011holomycin,lauinger2023} (Figure~1A), we removed DTT from the \textit{in vitro} transcription reaction and performed thiolutin treatment with or without molar equivalents of DTT (Figure~1B). Under the buffer conditions employed, we failed to observe thiolutin inhibition of Pol~II, with DTT having no effect (Figure~1B). We therefore set out to screen for potential cofactors.

\subsection{Three independent genetic screens for thiolutin-resistant and -sensitive mutants}

To gain insights into factors controlling the cellular response to thiolutin, we performed three independent genetic screens for thiolutin-resistant and -sensitive mutants in yeast. First, we conducted a forward genetic screen using UV mutagenesis followed by bulk segregant analysis and whole genome sequencing (Figure~S1A). Second, we screened a yeast Variomics library for individual thiolutin-resistant candidates \citep{pan2006}. Third, we performed high-throughput screens using both pooled Variomics and deletion libraries using barcode sequencing (Bar-seq) \citep{smith2009,lee2014} (Figure~2A).

Our data exhibited excellent reproducibility (Figure~S2) and revealed contributions of multidrug resistance (MDR) pathways, oxidative stress response (OSR) pathways, multiple divalent metal homeostasis/trafficking pathways, and the proteasome to thiolutin resistance.

\subsection{Multidrug resistance and oxidative stress pathways modulate thiolutin sensitivity}

From both manual Variomics and forward genetic screens, we isolated multiple mutants in \textit{YAP1} and \textit{YRR1}, encoding transcription factors implicated in OSR and/or MDR, respectively (Figure~S1C). We conclude that the MDR transcription factor Pdr1 is the major transcription factor promoting efflux of thiolutin from the cell, whereas neither Yrr1 nor Yap1 appeared to promote baseline resistance.

We also identified thiolutin-modulating mutations in OSR-related genes, including genes in the thioredoxin pathway. Notably, \textit{trr1$\Delta$} conferred hypersensitivity to thiolutin, while \textit{trx1$\Delta$ trx2$\Delta$} conferred thiolutin resistance and suppressed the \textit{trr1$\Delta$} hypersensitivity (Figure~S1G). These data indicate that accumulation of oxidized thioredoxins mediates thiolutin hypersensitivity and that Trx1 and Trx2 have a critical requirement for thiolutin activity \textit{in vivo}.

\subsection{Thiolutin induces apparent oxidative stress \textit{in vivo}}

Thiolutin induced Yap1::EGFP nuclear localization, consistent with induction of oxidative stress, although the thiolutin-induced nuclear translocation appeared slower and weaker than that induced by H$_2$O$_2$ (Figure~S5). Thiolutin depleted total glutathione within an hour of treatment (Figure~S6A), depleting both active GSH and GSSG. However, holomycin had a weak effect at best on Yap1 localization and did not affect cellular glutathione levels. We suggest that thiolutin is a redox cycler \textit{in vitro} and possibly also \textit{in vivo} but does not appear to inhibit transcription through widespread DNA damage.

\subsection{Thiolutin alters divalent metal homeostasis}

From our high-throughput screens, we observed distinct thiolutin resistance or sensitivity linked to diverse Zn$^{2+}$ trafficking genes (Figure~2B). The \textit{zap1$\Delta$} mutant conferred hypersensitivity to thiolutin, consistent with Zap1 function in cellular resistance to thiolutin and thiolutin alteration of Zn$^{2+}$ homeostasis \textit{in vivo}.

In addition to Zn$^{2+}$ homeostasis mutants, we observed genes involved in Mn$^{2+}$ trafficking linked to thiolutin resistance. The \textit{pmr1$\Delta$} mutant, which increases intracellular Mn$^{2+}$ levels, conferred thiolutin sensitivity (Figure~3B). Mn$^{2+}$ dose-dependently exacerbated thiolutin sensitivity in wild-type yeast (Figure~3C), supporting synergy between Mn$^{2+}$ and thiolutin \textit{in vivo}.

We confirmed that Zn$^{2+}$ caused a distinct UV shift of reduced thiolutin, and found that Mn$^{2+}$ and Cu$^{2+}$ can also interact with reduced but not non-reduced thiolutin in a manner likely similar to Zn$^{2+}$ binding (Figure~3E). NMR experiments confirmed these interactions (Figure~S9E).

\subsection{Thiolutin, when activated by DTT and Mn$^{2+}$, directly inhibits Pol~II \textit{in vitro}}

Remarkably, Mn$^{2+}$ addition to reduced thiolutin potently inhibited purified Pol~II activity (Figure~4A), in sharp contrast to the activation of Pol~II activity caused by Mn$^{2+}$ alone. Template binding protects Pol~II from the Thio/Mn$^{2+}$ inhibition (Figure~4B), suggesting that Thio/Mn$^{2+}$ inhibits Pol~II DNA interaction or a similarly early step in transcription, consistent with the behavior of a transcription initiation inhibitor.

We tested if Thio/Mn$^{2+}$ affected Pol~II elongation on a transcription bubble template using an RNA primer, where transcription initiation is bypassed (Figure~4E). Thio/Mn$^{2+}$ altered but did not block Pol~II transcription elongation, inducing a highly pause-prone elongation mode. Pre-assembled transcription elongation complexes were resistant to Thio/Mn$^{2+}$, validating the critical order-of-addition.

\subsection{Thiolutin reduces Pol~II occupancy genome-wide}

We observed a specific decrease of Pol~II occupancy by qPCR on the 5$'$ end of \textit{YLR454W} within the first 2 minutes of thiolutin treatment, consistent with fast transcription initiation inhibition (Figure~5A). ChIP-seq for FLAG-tagged Rpb3 showed broad decrease in Pol~II occupancy genome-wide (Figure~5C) with a continuum of effects where ribosomal protein (RP) genes are affected most strongly, then ribosome biogenesis (RiBi) genes, then other genes (Figure~5D).

\subsection{The apparent Thio/Mn$^{2+}$ complex is unstable in solution but Pol~II is stably altered}

We observed continuous changes in UV spectra after Mn$^{2+}$ was added to reduced thiolutin (Figure~6A), with the peak shifting from 340~nm to around 380~nm in the first two minutes, followed by decrease as a third species at 300~nm accumulated. Freshly prepared Thio/Mn$^{2+}$ strongly inhibited Pol~II but completely lost inhibitory activity after 20 minutes (Figure~6B). However, immediately treated Pol~II remained inhibited, suggesting stable alteration. Excess DTT abolished the inhibition, consistent with reductant-sensitive inhibition (Figure~6C).

\section{Discussion}

Our findings demonstrate that thiolutin directly inhibits RNA Polymerase~II transcription initiation \textit{in vitro} when activated by appropriate reduction and Mn$^{2+}$. This resolves a long-standing controversy about whether thiolutin can directly target Pol~II. The requirement for both reduction and Mn$^{2+}$ explains why many previous studies failed to observe direct inhibition---classical transcription buffers often contained Mn$^{2+}$ \citep{jimenez1973}, while most modern studies did not include Mn$^{2+}$ in their buffers.

The chemicogenetic screens reveal that thiolutin has complex effects \textit{in vivo} involving MDR efflux, oxidative stress through thioredoxin oxidation, and interactions with multiple divalent metals including Zn$^{2+}$, Mn$^{2+}$, and Cu$^{2+}$. The \textit{in vivo} effects on Pol~II occupancy are widespread but the major effects are consistent with Tor pathway inhibition and stress induction rather than simple direct Pol~II inhibition. This complexity argues strongly against the continued use of thiolutin as a simple transcription inhibitor for mRNA stability studies \citep{grigull2004,millan2006,pelechano2008}.

The unstable nature of the inhibitory Thio/Mn$^{2+}$ species, combined with its stable alteration of Pol~II, suggests either that Mn$^{2+}$ coordinates with reduced thiolutin to directly bind Pol~II at a site occluded by DNA, or that Mn$^{2+}$ facilitates formation of a covalent modification such as a disulfide bond within Pol~II. The reversibility by excess DTT is consistent with the latter model. Future structural studies will be needed to determine the precise mechanism of this inhibition.

\section{Materials and Methods}

\subsection{Yeast strains and reagents}
Yeast strains and primers used are listed in Table~S2. \textit{YAP1} C-terminal tagging with EGFP and gene deletion strains were constructed as described \citep{janke2004,longtine1998}. Chemicals were commercially obtained from the following: Cayman Chemical (Thiolutin), Gold Biotechnology (DTT, TCEP, 5FOA), Sigma (MnCl$_2$), JT Baker (ZnCl$_2$), BDH Chemicals (MgCl$_2$).

\subsection{Variomics screens}
Two separate pools of diploid Variomics libraries were used for genetic screens \citep{pan2006}. Haploid Variomics libraries were used for manual screening on SC + 10 $\mu$g/mL thiolutin plates. For high-throughput screening, pooled libraries were used with Bar-seq \citep{smith2009,lee2014}.

\subsection{\textit{In vitro} transcription assays}
Transcription activity assays were performed with ssDNA as the template and NTPs including $^{32}$P-labeled UTP. The transcribed $^{32}$P-containing RNA was separated from free $^{32}$P-UTPs on 10\% polyacrylamide gels for visualization. For elongation assays, transcription bubble templates with RNA primers were used to bypass initiation.

\subsection{ChIP-qPCR and ChIP-seq}
ChIP-qPCR was performed for FLAG-tagged Rpb3 subunit of Pol~II at a \textit{TEF1} promoter-driven \textit{YLR454W} reporter. For ChIP-seq, thiolutin (10 $\mu$g/mL) was added for 1, 2, 4, or 8 minutes, with DMSO vehicle control. The protocol included spike-in of Rpb3-FLAG chromatin from \textit{S. pombe} for normalization.

\subsection{UV spectroscopy and NMR}
UV spectra of reduced thiolutin treated with divalent metals were acquired at different time points. $^1$H-NMR experiments confirmed interactions between divalent metals and reduced thiolutin. Reduced thiolutin was prepared by treatment with equimolar DTT.

\subsection{Glutathione measurements}
Growing wild-type yeast cultures were washed and treated with indicated conditions at 30$^\circ$C for one hour. Reduced glutathione (GSH) and total glutathione were measured in three independent experiments.

\subsection{Statistical analysis}
Statistical analysis for \textit{in vitro} transcription experiments was performed using one-way ANOVA with multiple comparisons controlling for false discovery rate (0.01) using the method of Benjamini, Krieger, and Yekutieli \citep{benjamini1995}. For Yap1 localization, Friedman's paired test with Dunn's correction for multiple comparisons was used.

\bibliographystyle{plainnat}
\bibliography{references}

\end{document}
